\documentclass{article}
\usepackage{amssymb} % Required for \mathbb

\begin{document}

\section{Direct Proofs}

We are tying to prove 

$\forall x \in \mathbb{R}, p(x) \iff q(x)$

\\ \\ 

strategy:
1. assume the p(x) condition holds
2. Unwind the formula definion to get something we can use to do
calculations.
3. Show calculations lead to q(x)
4. repack an algorithm definition as a formal term to complete the process.

EX: if m and n are even ints, then m + n is also even.
How do we prove this to be generally true?

firstly, lets let p(x) be our first proposition ( m and n are even).
And q(x) be our second, which is m+n = even.

applying our strategy, we assume p(x) is true (m and n are assumed even then).

step two is to apply the formula/proof of an even number to m and n.
[rewrite m + n as the new definitions of their values.]

the sum then is 2(ksub1 + ksub2), which is even by the definition of an even number (which we assume is established).

{two questions: are we assuming the definition of an even number, or have we proved that elsewhere? secondly, the proof of closure is also assume.}

EX2 (direct proof of an odd integer)
if m is even and n is odd, then mn is even

again, assign p(x) and q(x) to our compound statements.

assume p(x) (m in even n is odd)



Take the defintion of an even number:
there exists a number k in Z such that 2k = n]

\section{More examples}

prove: if $n$ is odd, $n^2$ is odd

NOTE: for direct proofs, we want the first compount/p(x) to be 'simpler' than the second (q(x))

question:{can we only use contrapositive indirect proofs when the variables are binary?}

DEFN transitivity: A=B, B=C, then A=C

\end{document}
